\documentclass[10pt, xcolor=x11names, compress]{beamer}
%\documentclass[10pt, xcolor=x11names, compress, handout]{beamer}
\usetheme{progressbar}
%\usecolortheme[named=Purple4]{structure}
\progressbaroptions{headline=sections,titlepage=normal,frametitle=normal}

\setbeamertemplate{navigation symbols}{}

\usepackage{iwona} 

\usepackage{alltt}
\usepackage{amsmath,amsfonts, amssymb, amscd}
\usepackage{hyperref}
\usepackage{setspace}
\usepackage{wasysym}
\usepackage{ulem}

\usepackage{calc}
\usepackage[overlay,absolute]{textpos}
\TPGrid[5mm,5mm]{20}{20}



\renewcommand{\Re}{\operatorname{Re}}
\renewcommand{\Im}{\operatorname{Im}}
\newcommand{\debye}{\operatorname{debye}}

\newcommand{\chik}{$\chi(k)$}
\newcommand{\chir}{$|\tilde{\chi}(R)|$}


\newcommand{\file}[1]{{\color{Firebrick4}\texttt{`#1'}}}
\newcommand{\multiple}{{\color{Orange3}\textsl{multiple}}}


\newcommand{\atoms}  {{\color{DarkOrchid4}\textsc{atoms}}}
\newcommand{\feff}   {{\color{DarkOrchid4}\textsc{feff}}}
\newcommand{\ifeffit}{{\color{DarkOrchid4}\textsc{ifeffit}}}
\newcommand{\athena} {{\color{DarkOrchid4}\textsc{athena}}}
\newcommand{\artemis}{{\color{DarkOrchid4}\textsc{artemis}}}

\renewenvironment<>{center}
{\begin{actionenv}#1\begin{originalcenter}}
{\end{originalcenter}\end{actionenv}}

\definecolor{guessp}   {rgb}{0.64,0.00,0.64}
\newcommand{\guessp}   {{\color{guessp}guess}}
\definecolor{defp}     {rgb}{0.00,0.55,0.00}
\newcommand{\defp}     {{\color{defp}def}}
\definecolor{setp}     {rgb}{0,0,0}
\newcommand{\setp}     {{\color{setp}set}}
\definecolor{lguessp}  {rgb}{0.24,0.11,0.56}
\newcommand{\lguessp}  {{\color{lguessp}lguess}}
\definecolor{skipp}    {rgb}{0.70,0.70,0.70}
\newcommand{\skipp}    {{\color{skipp}skip}}
\definecolor{restrainp}{rgb}{0.80,0.61,0.11}
\newcommand{\restrainp}{{\color{restrainp}restrain}}
\definecolor{afterp}   {rgb}{0.29,0.44,0.55}
\newcommand{\afterp}   {{\color{afterp}after}}
\definecolor{penaltyp} {rgb}{0.55,0.35,0.17}
\newcommand{\penaltyp} {{\color{penaltyp}penalty}}
\definecolor{mergep}   {rgb}{0.93,0.00,0.00}
\newcommand{\mergep}   {{\color{mergep}merge}}


%% define new commands here
%\newcommand{\eto}{EuTiO$_3$}

\mode<presentation>

\title{Mechanics of an XAS Scan}
%\subtitle{}

\author{Bruce Ravel}
\institute[NIST]{Synchrotron Methods Group, Materials Measurement Science Division\\%
  Materials Measurement Laboratory\\%
  National Institute of Standards and Technology\\%
  \&\\%
  Local Contact, Beamline X23A2\\%
  National Synchrotron Light Source\\~}


%\date[Diamond2011]{EXAFS Data Analysis workshop 2011\\
  Diamond Light Source\\November 14--17, 2011\\~}

\date[\today]{\today}

\begin{document}
\maketitle

\begin{frame}
  \frametitle{Copyright}
  \tiny

  This document is copyright \copyright 2007-2010 Bruce Ravel.

  \begin{center}
    \includegraphics[width=1.0cm]{images/somerights20}
  \end{center}

  This work is licensed under the Creative Commons
  Attribution-ShareAlike License.  To view a copy of this license,
  visit \href{http://creativecommons.org/licenses/by-sa/3.0/}
  {\color{Purple4}\texttt{http://creativecommons.org/licenses/by-sa/3.0/}}
  or send a letter to Creative Commons, 559 Nathan Abbott Way,
  Stanford, California 94305, USA.

  \begin{description}
  \item[You are free:] %
    \begin{itemize}
    \item \textbf{to Share} --- to copy, distribute, and transmit the work
    \item \textbf{to Remix} --- to adapt the work
    \end{itemize}
  \item[Under the following conditions:] %
    \begin{itemize}
    \item Attribution. You must attribute the work in the manner
      specified by the author or licensor (but not in any way that
      suggests that they endorse you or your use of the work).
    \item Share Alike. If you alter, transform, or build upon this
      work, you may distribute the resulting work only under the same,
      similar or a compatible license.
    \item Any of these conditions can be waived if you get permission
      from the author.
    \end{itemize}
  \end{description}
  \begin{itemize}
  \item For any reuse or distribution, you must make clear to others
    the license terms of this work. The best way to do this is with a
    link to the URL for this document.
  \item Any of the above conditions can be waived if you get
    permission from the copyright holder.
  \item Nothing in this license impairs or restricts the author's
    moral rights.
  \end{itemize}

  Your fair dealing and other rights are in no way affected by the
  above.  This is a human-readable summary of the Legal Code (the full
  license).


\end{frame}

%%% Local Variables:
%%% mode: latex
%%% TeX-master: "pimst2"
%%% End:


\begin{frame}
  \frametitle{Good data starts with a good plan}
  
\end{frame}

\begin{frame}
  \frametitle{Two scan types}
  \begin{description}
  \item[Step scan] blah blah
  \item[Slew scan] Blah blah
  \end{description}
\end{frame}

\section{Step scanning}

\begin{frame}
  \frametitle{XANES and EXAFS}
  
\end{frame}

\begin{frame}
  \frametitle{Scan parameters}
  \begin{enumerate}
  \item Regions
  \item Step sizes
  \item Integration times
  \end{enumerate}
\end{frame}

\begin{frame}
  \frametitle{Considerations for data processing}
  Show a figure of canonical XAS data (some Fe oxide) with pre- and
  post edge lines.  Overplot with a severely truncated version of the
  same.  Discuss the difficulties of normalization and its impact on
  subsequent interpretation.
\end{frame}

\section{Slew scanning}

\begin{frame}
  \frametitle{Slew...}
  
\end{frame}





\end{document}

%%% Local Variables:
%%% mode: latex
%%% TeX-master: t
%%% TeX-parse-self: t
%%% TeX-auto-save: t
%%% TeX-auto-untabify: t
%%% TeX-PDF-mode: t
%%% End:
\grid
